\section{R-NEW-004: the hook-chain Lie stabilizer is zero}

Let $W=\mathbb C^n$ with ordered basis $e_1,\ldots,e_n$.
Realize $\mathbb S_{(2,1)}W$ as the degree-three free Lie
submodule of $W^{\otimes3}$, with bracket $[a,b]=ab-ba$.
For clarity, the map
$W\otimes\Lambda^2W\to W^{\otimes3}$ taking
$a\otimes(b\wedge c)$ to $[a,[b,c]]$ has kernel
$\Lambda^3W$, by Jacobi.  This can also be checked separately
in each multidegree: repeated-letter multidegrees have image
dimension one and distinct-letter multidegrees have image
dimension two.  Thus its image is the usual Schur module
$\mathbb S_{(2,1)}W$; the identification is equivariant.

\begin{lemma}[R-NEW-004]
For every $n\ge2$, the vector
\[
 z_n=(e_n,F_n)\in W\oplus\mathbb S_{(2,1)}W,
 \qquad F_n=\sum_{i=1}^{n-1}[e_i,[e_i,e_{i+1}]]
\]
has zero $\mathfrak{gl}_n$ stabilizer.  Consequently its
algebraic $GL_n$ stabilizer over characteristic zero is finite.
\end{lemma}

\begin{proof}
Write $X e_j=\sum_a x_{aj}e_a$ and suppose $Xz_n=0$.
The first summand gives $Xe_n=0$.  The action on $F_n$ is
the derivation action on the displayed brackets.

For $r\ne s$, the repeated-letter multidegree
$2e_r+e_s$ is spanned by the nonzero vector
$L_{rs}=[e_r,[e_r,e_s]]$.  Its coefficient in $XF_n$ is
\[
 {\bf1}_{r<n}x_{s,r+1}
       +2{\bf1}_{s=r+1}x_{rr}
       -{\bf1}_{s=r-1}x_{r,r-1}.                 \tag{R4-1}
\]
The last term follows from
$[e_r,[e_{r-1},e_r]]=-L_{r,r-1}$; the other possible
derivative of that repeated letter is zero.  Formula(R4-1)
accounts for every term of that multidegree.

Taking $s=r+1$ gives
$2x_{rr}+x_{r+1,r+1}=0$.  Since $x_{nn}=0$, all diagonal
entries vanish.  For column $j\ge2$, (R4-1) then says that
the only possibly nonzero entries are
\[
 x_{j-1,j},\qquad
 x_{j-2,j}=x_{j-1,j-2}\quad(j\ge3).
                                                        \tag{R4-2}
\]
In particular every entry strictly below the diagonal in
columns $j\ge2$ is zero.  For $j\ge4$, the right side of
the second equality in(R4-2) is such an entry and is zero.
Thus, besides the first column and the first superdiagonal,
the sole possible entry is $x_{13}=x_{21}$ when $n\ge3$.
Put $y_j=x_{j-1,j}$; the equation $Xe_n=0$ gives $y_n=0$.

If $n=2$, the same equation kills the remaining upper entry,
and (R4-1) at $(r,s)=(2,1)$ gives $x_{21}=0$.  Hence assume
$n\ge3$.

For each $a\ge4$, the distinct-letter multidegree
$e_1+e_2+e_a$ receives a contribution only from the term
with index1, with coefficient $x_{a1}$.  Its Lie vector
$[e_a,[e_1,e_2]]+[e_1,[e_a,e_2]]$ is nonzero, as is seen
by its associative word expansion.  Hence $x_{a1}=0$.

For $3\le j\le n-1$, the multidegree
$e_{j-1}+e_j+e_{j+1}$ can come from only the two adjacent
terms of $F_n$.  Their coefficients are $y_j$ and
$x_{j+1,j-1}$ respectively.  The latter is already zero
by(R4-2), while the former multiplies the nonzero vector
\[
 [e_{j-1},[e_j,e_{j+1}]]+
 [e_j,[e_{j-1},e_{j+1}]].
\]
Therefore every such $y_j$ is zero.

At multidegree $e_1+e_2+e_3$, take the independent basis
\[
 A=[e_1,[e_2,e_3]],\qquad
 B=[e_2,[e_1,e_3]].
\]
Jacobi gives $[e_3,[e_1,e_2]]=-A+B$.
The remaining contribution is exactly
\[
                       y_2(A+B)+x_{31}(-2A+B).
\]
Thus $y_2-2x_{31}=y_2+x_{31}=0$, and characteristic zero
implies $x_{31}=y_2=0$.

It remains to kill $x_{13}=x_{21}$.  For $n=3$ this follows
already from $Xe_3=0$.  For $n\ge4$, the multidegree
$e_1+e_3+e_4$ can only receive the derivative of the
repeated $e_3$ in the index3 term.  It has coefficient
$x_{13}$ times a nonzero Lie vector, so $x_{13}=0$.
Every entry of $X$ has now vanished.  An algebraic
stabilizer in characteristic zero has dimension equal to
the dimension of its Lie algebra.  It is therefore
zero-dimensional and finite, proving the last assertion.
\end{proof}

The proof is uniform in $n$ and uses no finite-rank extrapolation.
The separate exact script reconstructs the associative-word
action matrices and checks full rank $n^2$ in small dimensions
as a convention check.  No modular characteristic-three Lie
smoothness is claimed: the $2\times2$ coefficient system in
the distinct-letter step has determinant three.
